Der Hammer balanciert in seinem Schwerpunkt: Dort heben sich die Drehmomente der beiden Gewichtskräfte auf. Sei $\ell$ die Länge des Stiels, $m_1$ seine Masse und $m_2$ die Masse des Kopfs. Der Schwerpunkt des Stiels liegt $\ell/2$ vom Griffende entfernt, der des Kopfs $\ell$.

Für den Abstand $x$ des Balancierpunkts vom Griffende gilt (Drehmomente bezüglich des Balancierpunkts)
\begin{align}
 m_1\,g\cdot\left(x-\frac{\ell}{2}\right) &= m_2\,g\cdot(\ell-x)
\end{align}
und damit
\begin{align}
 x &= \xF \\
   &= \frac{\mH\cdot\frac{\LS}{2}+\mK\cdot\LS}{\mH+\mK} \\
   &= \x \approx \result{\xP{-2}{2}}
\end{align}
Nahe beim Kopf: Der schwere Kopf zieht den Schwerpunkt zu sich.
